\label{section:delay}
In the Delay model the Algorithm ca wait a certain amount of time before matching a vertex

\begin{theorem}
	For delay no greater than $n/2$ any deterministic online algorithm has competitiveness of $0.5$
\end{theorem}

\begin{proof}
	Consider the following graph $G$: There are $k=n/2$ groups $\{Z_i|i\in[k]\}$. of two offline vertices\\
	First, $k$ vertices $z_i^D$ arrive, each connected to both vertices of $Z_i$.\\
	As $z_k^D$ arrives the cache is full, and $z_1^D$ has to be matched. Let the offline vertex it was matched to be $z_1^B$. Next, $z_1^C$ arrives, only connected to $z_1^B$, which is unmatchable. Let $z_1^A$ be the unmatched offline vertex.\\
	After $z_1^A$ has arrived $z_2^D$ has to be matched. The procedure above is repeated until $z_k^C$ has arrived.\\
	The perfect matching of $G$ connects all $z_i^A,z_i^D$ and $z_i^B,z_i^C$ pairs, while the algorithm connects only $z_i^B,z_i^D$ pairs, making the competitiveness $0.5$. The final graph is illustrated in the figure below.
\end{proof}

\begin{figure}[H]
	\centering
	\scalebox{.8}{\begin{tikzpicture}
	\begin{pgfonlayer}{nodelayer}
		\node [style=default] (0) at (-1, 4) {};
		\node [style=default] (1) at (1, 4) {};
		\node [style=default] (2) at (-1, 3) {};
		\node [style=default] (3) at (1, 3) {};
		\node [style=none] (4) at (-1, 4.5) {$z_1^A$};
		\node [style=none] (5) at (1, 4.5) {$z_1^D$};
		\node [style=none] (6) at (-1, 3.5) {$z_1^B$};
		\node [style=none] (7) at (1, 3.5) {$z_1^C$};
		\node [style=none] (8) at (-1.5, 4) {};
		\node [style=none] (9) at (-1.5, 3) {};
		\node [style=none] (10) at (-2, 3.5) {$Z_1$};
		\node [style=default] (11) at (-1, 2) {};
		\node [style=default] (12) at (1, 2) {};
		\node [style=default] (13) at (-1, 1) {};
		\node [style=default] (14) at (1, 1) {};
		\node [style=none] (15) at (-1.5, 2) {};
		\node [style=none] (16) at (-1.5, 1) {};
		\node [style=default] (17) at (-1, -0.25) {};
		\node [style=default] (18) at (1, -0.25) {};
		\node [style=default] (19) at (-1, -1.25) {};
		\node [style=default] (20) at (1, -1.25) {};
		\node [style=none] (21) at (-1.5, -0.25) {};
		\node [style=none] (22) at (-1.5, -1.25) {};
		\node [style=none] (23) at (-2, 1.5) {$Z_2$};
		\node [style=none] (24) at (-2, -0.75) {$Z_k$};
		\node [style=none] (25) at (0, 0.75) {.};
		\node [style=none] (26) at (0, 0.5) {.};
		\node [style=none] (27) at (0, 0.25) {.};
	\end{pgfonlayer}
	\begin{pgfonlayer}{edgelayer}
		\draw (0) to (1);
		\draw (2) to (1);
		\draw (2) to (3);
		\draw [bend right=90, looseness=0.75] (8.center) to (9.center);
		\draw (11) to (12);
		\draw (13) to (12);
		\draw (13) to (14);
		\draw [bend right=90, looseness=0.75] (15.center) to (16.center);
		\draw (17) to (18);
		\draw (19) to (18);
		\draw (19) to (20);
		\draw [bend right=90, looseness=0.75] (21.center) to (22.center);
	\end{pgfonlayer}
\end{tikzpicture}
}
	\label{fig:det-delay-figure}
	\caption{Worst case scenario for a deterministic cache algorithm}
\end{figure}

\begin{theorem}
	For any $f(n)=o(n)$ delay, no algorithm has asymptotic competitiveness better than $1-\frac{1}{e}$
\end{theorem}
\begin{proof}
	For $\mathrm{ALG}$ construct the adversary graph in a manner similar to the one above, with $f(n)$ copies of $G_A$, a graph on which $\mathrm{ALG}$ has a performance of $1-\frac{1}{e}$. Notice that the size of individual $G_A$ is $\frac{n}{f(n)}$ and since $\lim_{n\to\infty}\frac{n}{f(n)}\to\infty$, the size of individual $G_A$ instances is growing with $n$, leading to $\mathrm{ALG}$ being forced to solve isolated instances of size growing to infinity, yielding the stated competitiveness.
\end{proof}